交通预测系统能够获得的历史观测往往不完整，而且在不断变化：新部署的传感器逐步积累读数，在役传感器可能中断，部分位置在训练期没有任何动态信号。本文研究路网拓扑已知、部分节点不参与训练时的动态观测支持下的归纳时空预测，将窗口内可用观测的位置、数量与时间分布（观测支持）既视为输入有效性的描述，也视为组织预测计算的条件。为此我们提出 ALFNet。它以逐步的观测角色、状态年龄和观测掩码构成生命周期上下文；共享的图循环补全先为两个分支提供同一份历史；图分支依据观测状态对五种传播算子加权组合，并以残差自回归方式解码；时间分支从空间分区上下文生成条件特征，通过 FiLM 调制图编码状态；输出端再施加由末步角色控制幅度的有界图修正。在 METR-LA、PEMS-BAY、PEMS04 和 PEMS08 上，我们以零历史（Unseen）、持续观测（Steady）、逐步加入（Join）和历史中断（Fault）四类条件统一评价。与所列八个比较方法（共同协议下的适配版本）相比，ALFNet 的跨条件平均误差（LifeAvg）三项指标在四个数据集上均为最优，相对各数据集最强基线的 MAE 降低7.86\%–15.63\%；逐步加入（Join）与历史中断（Fault）两类主条件的三项指标在四个数据集上均为最优。组件消融表明跨节点时间上下文的贡献最大，观测状态条件化传播与状态年龄带来较小但方向基本一致的改进。收益主要出现在本地历史部分存在的条件；零历史（Unseen）条件下的优势依赖数据集，在两个流量数据集上未能成立。
